\section{Reproducibility}
\label{sec:repro}

This section is what somebody would need to rerun the work and get the same
answers, and an explicit statement of which numbers in this build of the report
were produced under the current protocol and which were not.

\subsection{The machine and the toolchain}

\begin{itemize}
\item GPU: NVIDIA RTX 5070, Blackwell GB205, compute capability 12.0, sm\_120, 48
      SMs, 12 GB GDDR7.
\item Driver: 610.62, reached from WSL2 through the Windows driver, which is why
      the profiler counter permission is a Windows side setting and not a Linux
      one.
\item Toolkit: CUDA 13.3, nvcc 13.3.73. The CMake floor is toolkit 12.8, the first
      release with sm\_120 codegen, guarded with a message rather than a link
      error.
\item Host compiler: g++-14.3, and this one is not optional. The system default on
      this machine is GCC 15.2, whose libstdc++ headers use \texttt{if consteval}
      in a form that the nvcc 13.3 frontend cannot parse, so a default configure
      fails inside a system header with no obvious connection to this project. The
      build pins the host compiler explicitly; \texttt{docs/building.md} carries the
      diagnosis and the exact CMake flags.
\item Build: CMake 4.4.0 with Ninja, release architecture list
      \texttt{80-real;86-real;89-real;90-real;120-real;120-virtual}. Warnings are
      errors.
\item Report toolchain: matplotlib 3.10.7 and numpy 2.3.5, pinned exactly rather
      than as floors, because matplotlib writes its own rasteriser output into the
      PNG and its bundled FreeType into the PDF, so two different wheels draw two
      different files from the same numbers. The report target exports
      \texttt{SOURCE\_DATE\_EPOCH} to a fixed instant so the PDF timestamp is not a
      moving part either, and CI runs
      \texttt{git diff -{}-exit-code report/figures report/tables} after rebuilding
      them.
\end{itemize}

\subsection{The measurement protocol}

\begin{itemize}
\item \textbf{Clocks.} The 1.1.0 sweep locks the graphics clock before it measures
      anything and refuses to run if it cannot. \texttt{docs/benchmarking.md} says
      what to do when the clock cannot be locked, and the answer is not to measure
      anyway.
\item \textbf{Timing.} CUDA events around the launch, warmup launches discarded,
      median of the timed repetitions reported with the interquartile range beside
      it.
\item \textbf{Independent runs.} Five independent process repeats per
      configuration under the 1.1.0 protocol, reduced to one row of record with a
      bootstrap confidence interval carried in the row. The v1 sweep was a single
      process with five warmup and twenty timed repetitions, and carries no
      interval.
\item \textbf{Order and thermal state.} Configurations run in shuffled order with
      cooldowns between them. NVML samples power, clocks, temperature and throttle
      reasons throughout; a throttled configuration fails the run rather than being
      reported.
\item \textbf{Baselines.} Vendor baselines are measured once per shape, in the same
      process, on the same device, with no setup work inside the timed region. An
      allocation gate arms a CUPTI callback on the allocation entry points around
      each timed lambda and fails the run if the counter moves.
\item \textbf{cuBLAS math mode.} \texttt{bench\_all} sets the math mode explicitly,
      reads it back, and records the result in the \texttt{cublas\_math\_mode}
      column of every row. The default is \texttt{CUBLAS\_DEFAULT\_MATH}, which
      permits reduced precision reduction of split-K partials; the
      \texttt{-{}-disallow-reduced-precision-reduction} flag adds
      \texttt{CUBLAS\_MATH\_DISALLOW\_REDUCED\_PRECISION\_REDUCTION} so the effect
      on both speed and oracle accuracy can be measured rather than assumed.
\item \textbf{Workspace policy.} \texttt{cublasSetWorkspace} is never called, at
      any size. One call forfeits the default pool, which is 32 MiB on this part,
      and a baseline running without its pool is not the baseline a user would get.
      \texttt{cublasSetStream} is called first, because it resets the workspace to
      the default pool. On my own side, a Context that owns a workspace keeps
      split-K and stream-K scratch allocation out of the timed call, and
      \texttt{ckl::gemm\_workspace\_size} is how a caller sizes it.
\item \textbf{Provenance.} Every results row records the commit that produced it,
      and \texttt{make check-style} runs a gate that resolves each recorded hash in
      git and fails when one does not.
\end{itemize}

\subsection{The exact targets}

\begin{verbatim}
make setup      # configure (CMake plus Ninja, sm_120)
make build      # compile; zero warnings is a gate
make test       # correctness against cuBLAS, cuSPARSE, cuFFT, CUB, CPU
make sweep      # full protocol sweep -> experiments/results/summary.csv
make roofline   # measure the ceilings and render the roofline figure
make report     # regenerate figures and tables, then build both PDFs
make check-style
make all        # build, test, sweep, report, style gate
\end{verbatim}

\texttt{make report} does not need a GPU or a toolkit. It works from the committed
results files, which is the point: a clean clone on a machine with no NVIDIA
hardware rebuilds every figure and table in this document, and fails loudly if an
input is missing.

\subsection{Which numbers here are which}

This matters more than the rest of the section, so it is stated bluntly.

\textbf{Every performance number in this build of the report is v1 era.} The
committed \texttt{experiments/results/summary.csv} is the v1 sweep: single process,
unlocked clocks, and rows whose commit hash no longer resolves. That covers the
GEMM ladder table, the ladder figure, the gap decomposition, the supporting families
table, the roofline points, and every percent of a vendor library quoted anywhere in
this document. Read them as provisional. The withdrawn cuSPARSE percentages are not
merely provisional; they are gone, for the reason in Section \ref{sec:limitations}.

\textbf{Every compile time fact here is current.} The register allocation study, the
SASS instruction mix, the tile family's register and shared memory footprints, and
the shape constraints are produced by the 1.1.0 toolchain from the current tree. They
need no GPU, so they do not depend on the clock protocol at all.

\textbf{Every 1.1.0 family performance number reads pending.} The tile family, the
predicated tails, split-K, stream-K, the CUTLASS reference line, the SpMV matrix
suite, FFT and convolution, and reduction and scan have no throughput rows in the
committed summary, so their tables carry the word \emph{pending} in place of a
number. The generator writes it there deliberately, and no estimate is substituted.
The owner sweep at locked clocks is what fills them in, and it is the one action
that closes both this section and the corrections register.
